\documentclass[11pt]{article}
\usepackage[T1]{fontenc}
\usepackage{textcomp}
\usepackage[margin=0.75in]{geometry}
\usepackage{amsmath,amssymb,amsthm}
\usepackage{array,booktabs}
\usepackage{hyperref}
\setlength{\parindent}{0pt}
\newtheorem{theorem}{Theorem}
\theoremstyle{definition}
\newtheorem{definition}{Definition}
\title{Flow Proof Artifact: prop-01-equal-chords-equal-angles}
\author{Generated by \texttt{flow doc proof}}
\date{Flow Proof Book}
\begin{document}
\maketitle
In equal circles equal chords subtend equal angles at the center, and equal central angles subtend equal chords.\\[0.5em]
\textbf{Source.} euclid — Elements, Book III, Proposition 1\\[0.5em]
\bigskip
\noindent{\large \textbf{Derived fact 1.} \textit{In equal circles equal chords subtend equal angles at the center, and equal central angles subtend equal chords} \hfill the Euclidean plane \cdot Euclid Book III \cdot Proposition 1: equal chords subtend equal angles at the center}\par
\smallskip\noindent\textit{Coordinate: the Euclidean plane \cdot Euclid Book III \cdot Proposition 1: equal chords subtend equal angles at the center}\par
\smallskip\noindent\textit{Source: euclid — Elements, Book III, Proposition 1}\par
\medskip
\noindent\textbf{Goal.} In equal circles equal chords subtend equal angles at the center, and equal central angles subtend equal chords\par
\noindent$\displaystyle \text{chord AB} = \text{chord DE} implies \angle AOB = \angle DOE$\par
\medskip
\noindent
\begin{tabular}{@{} >{\raggedright\arraybackslash}p{0.06\textwidth} >{\raggedright\arraybackslash}p{0.47\textwidth} >{\raggedleft\arraybackslash}p{0.41\textwidth} @{}}
\textbf{\#} & \textbf{Proof} & \textbf{Mathematics} \\
\midrule
\textbf{1.} & We prove that proposition 1: equal chords subtend equal angles at the center for Euclid Book III the Euclidean plane. & \\[0.45em]
\textbf{2.} & Let O = center\_of\_circle. & \\[0.45em]
\textbf{3.} & Let AOB = angle\_at\_center. & \\[0.45em]
\textbf{4.} & Let DE = equal\_chord. & \\[0.45em]
\textbf{5.} & From step 2, step 3, and step 4, we can deduce that angle AOB equals angle DOE. & $\displaystyle \angle AOB = \angle DOE$ \\[0.45em]
\textbf{6.} & We can deduce that chord AB equals chord DE implies angle AOB equals angle DOE. Hence proven. & $\displaystyle \text{chord AB} = \text{chord DE} implies \angle AOB = \angle DOE$ \\[0.45em]
\bottomrule
\end{tabular}
\label{thm:Geometry-Euclid Book III-proposition_1_equal_chords_subtend_equal_angles_at_the_center}
\par\medskip\hrule
\end{document}
